\documentclass{exam}

\newif\ifA
\newif\ifkey
\newif\ifprintselected   % required by exam_config's \ansbox

%%%%%%%%%%%%%%%%%%%%%%%%%%%%%%%%%%%%%%%%%%%%%%%%%%%%%%
%%% SET THESE IF VARIABLES
%%%%%%%%%%%%%%%%%%%%%%%%%%%%%%%%%%%%%%%%%%%%%%%%%%%%%%
% Lab Num -- Module 3's six lab periods follow L04 (Lab 1) and the two
% Module 2 measurement labs, so L21 is Lab 6. Override with \def\labNumIn{}
% if the course sequence renumbers.
\ifdefined\labNumIn
	\def\labNum{\labNumIn}
\else
	\def\labNum{6}
\fi

% is this the key or not?
\ifdefined\iskey
	\ifnum\iskey=1
		\keytrue
	\else
		\keyfalse
	\fi
\else
	%% Pick one manually
	\keyfalse % if this is NOT the key
%	\keytrue     % if this IS the key
\fi

% set up the header and footer - change for every term
\ifdefined\thisTermIn
	\def\thisterm{\thisTermIn}
\else
	\def\thisterm{Fall 2026}
\fi
%
\def\thisclass{ECE 444}

%%%%%%%%%%%%%%%%%%%%%%%%%%%%%%%%%%%%%%%%%%%%%%%%%%%%%%
%%% END SET THESE IF VARIABLES
%%%%%%%%%%%%%%%%%%%%%%%%%%%%%%%%%%%%%%%%%%%%%%%%%%%%%%
% packages and shortcuts
\input{myPackages_gr}
\usepackage[hidelinks, linkcolor=red]{hyperref}
\input{myShortcuts}
\setlength{\parindent}{0pt}
\setlength{\parskip}{6pt}

%% exam class customization
\input{exam_config}
\printselectedfalse      % full worked solutions in the key, not answers-only
\CorrectChoiceEmphasis{\color{red}}
\SolutionEmphasis{\color{red}}
\renewcommand\questionlabel{\textbf{\thequestion.}}
\renewcommand{\thepartno}{\alph{partno}}
\renewcommand{\partlabel}{(\thepartno)}

% Fillable table cell: blank in the student copy, the expected value in red on
% the KEY. Fixed-width centered columns keep the blank cells writable.
\newcolumntype{K}[1]{>{\centering\arraybackslash}p{#1}}
\newcommand{\kv}[1]{\ifkey{\color{red}#1}\fi\rule[-0.10in]{0pt}{0.34in}}

% print the answers if this is the key
\ifkey
	\printanswers % if this IS the key
\else
	\noprintanswers  %if this is NOT the key
\fi

% print the header and footer
% need to print key on top if this is the key
\ifkey
	\firstpageheader{}{\color{red}{ KEY}}{}
	\runningheader{\thisclass\hspace{1pt} -- \thisterm }{Lab \#\labNum\hspace{1pt}\color{red}{ KEY}}{\thepage\hspace{1pt} of\hspace{1pt} \numpages}
	\firstpagefooter{}{}{}
	\runningfooter{}{}{}
\else
	\firstpageheader{}{}{}
	\runningheader{\thisclass \hspace{1pt} -- \thisterm }{Lab \#\labNum\hspace{1pt}}{\thepage\hspace{1pt} of\hspace{1pt} \numpages}
	\firstpagefooter{}{}{}
	\runningfooter{}{}{}
\fi

\runningheadrule

\begin{document}
\pagestyle{headandfoot}   % exam class


%%%%%%%%%%%%%%%%%%%%%%%%%%%%%%%%%%%%%%%%%%%%%%%%%%%%%%
% title block
%%%%%%%%%%%%%%%%%%%%%%%%%%%%%%%%%%%%%%%%%%%%%%%%%%%%%%

\begin{center}
USAF ACADEMY \\
DEPARTMENT OF ELECTRICAL AND COMPUTER ENGINEERING \\
LAB \#\labNum\ (Lesson 21) - Array Factor Lab \\
\textbf{Lab Sheet} \\
\thisterm \\[4mm]
\end{center}

\vspace{-4pt}
\textbf{Name:} \rule[-2pt]{2.6in}{0.4pt} \hfill
\textbf{Section:} \rule[-2pt]{0.9in}{0.4pt} \hfill
\textbf{Date:} \rule[-2pt]{1.1in}{0.4pt}

\vspace{4pt}

\fullwidth{\textbf{Learning Objective 3.2: } Derive and apply the array factor for a
uniform linear array, and relate its beamwidth and sidelobe structure to what a
beam sweep measures.}

\vspace{2pt}

This sheet is the turn-in document for the Lesson 21 lab. Fill the calculated
columns before you sweep and the measured columns at the bench, so the
prediction is on paper before the trace is on the screen. The lab is a team
assignment, so one sheet per team.

\begin{questions}

%%%%%%%%%%%%%%%%%%%%%%%%%%%%%%%%%%%%%%%%%%%%%%%%%%%
\question \textbf{The measurement table.} One row per aperture, at
$10.3\ \ghz$ with $d = 14$ mm. Report each trace's peak level, and the step from the
eight-element peak beside it, and report each sidelobe level relative to that trace's own
peak.

\begin{center}
\small
\setlength{\tabcolsep}{4pt}
\renewcommand{\arraystretch}{1.25}
\begin{tabular}{| l | K{0.72in} | K{0.62in} | K{0.62in} | K{0.85in} | K{0.62in} | K{0.85in} |}
\hline
\textbf{Active elements} & \textbf{Peak (dBFS)} & \textbf{HPBW meas} & \textbf{HPBW calc} & \textbf{FNBW meas} & \textbf{FNBW calc} & \textbf{First SLL} \\\hline\hline
8 (Rx1-Rx8) & \kv{reference} & \kv{$13.1\mydeg$} & \kv{$13\mydeg$} & \kv{$28.1\mydeg$} & \kv{$30.1\mydeg$} & \kv{$-11$ to $-13$ dBc} \\\hline
4 (Rx3-Rx6) & \kv{$-6.0$ dB rel.} & \kv{$29\mydeg$} & \kv{$27\mydeg$} & \kv{$60$ to $62\mydeg$} & \kv{$62\mydeg$} & \kv{$-10$ to $-12$ dBc} \\\hline
2 (Rx4, Rx5) & \kv{$-11.5$ dB rel.} & \kv{$65\mydeg$} & \kv{$62\mydeg$} & \kv{no null in the scan} & \kv{$180\mydeg$} & \kv{none} \\\hline
\end{tabular}
\end{center}

\begin{solution}
The calculated column uses $\theta_{\text{HP}} \approx 0.886\lambda/(Nd)$ and
$\text{FNBW} = 2\arcsin(\lambda/Nd)$ with $\lambda = 29.1$ mm and $d = 14$ mm.
The two-element entries are the exceptions: the $62\mydeg$ comes from solving
$\cos(\psi/2) = 1/\sqrt{2}$ directly, and the $180\mydeg$ is a convention,
because a null there needs $\sin\theta = \lambda/2d = 1.04$ and no real angle
satisfies that.

Three measured deviations are expected, and each has a number attached. The
eight-element FNBW reads $28.1\mydeg$ rather than $30.1\mydeg$ because the true
null at $15.1\mydeg$ falls between grid samples and the nearest sample sits at
$14.06\mydeg$. The first step in peak level is very close to the $6.0$ dB the
voltage sum predicts, and the second is nearer $5.5$ dB, because with only two
elements on the peak sits about 11 dB above the floor and the reported power
includes the noise, which lifts the reading by roughly $0.3$ dB. The
four-element sidelobe reads high for the same reason, since the floor is then
only about 17 dB below that trace's peak. Sidelobe entries are reported as a
range and not as a single number at this dynamic range. In simulation the three
sweeps read $13.1\mydeg$, $29.1\mydeg$, and $65.4\mydeg$.
\end{solution}

%%%%%%%%%%%%%%%%%%%%%%%%%%%%%%%%%%%%%%%%%%%%%%%%%%%
\question \textbf{Written answer 1.} Explain why the first-null beamwidth of the
two-element array is quoted as $180\mydeg$ rather than measured. Name the
condition on $\lambda/2d$ that decides whether a null exists in visible space,
and evaluate it for $d = 14$ mm at $10.3\ \ghz$.

\begin{solutionorlines}[2.6in]
For two elements the array factor is $\cos(\psi/2)$ with $\psi = kd\sin\theta$,
so the first null needs $\psi = \pi$, which is $\sin\theta = \lambda/2d$. A null
exists in visible space only when $\lambda/2d \le 1$, that is when
$d \ge \lambda/2$. At $10.3\ \ghz$ the wavelength is $29.1$ mm, so
\eq{ \frac{\lambda}{2d} = \frac{29.1}{2(14)} = \frac{29.1}{28} = 1.04 > 1 }
and no real angle satisfies the condition. The null exists in the mathematics
but falls outside the range of angles the array can see, so the pattern has no
null anywhere in the scan and the beamwidth between first nulls is quoted as the
full $180\mydeg$ of the scan range.
\end{solutionorlines}

%%%%%%%%%%%%%%%%%%%%%%%%%%%%%%%%%%%%%%%%%%%%%%%%%%%
\question \textbf{Written answer 2.} The gain of each patch element is the same
whether the other seven elements are on or off, and the amplifier behind each
element is unchanged. Explain why the peak of the sweep nonetheless drops by
about 6 dB every time you halve the number of active elements.

\begin{solutionorlines}[2.4in]
At boresight the element signals arrive in phase, so the beamformer adds them as
voltages rather than as powers. Eight equal in-phase signals give eight times the
voltage of one, four give four times, and the plotted peak follows
\eq{ 20\mylog{\frac{N_{\text{active}}}{8}} = -6.0 \text{ dB at } N = 4,
\qquad -12.0 \text{ dB at } N = 2 }
Nothing about the individual element changed. What changed is how many copies of
the same signal are being summed, and the summed voltage halves each time the
active count halves. The measured second step comes out nearer $5.5$ dB, because
by then the peak is only about 11 dB above the noise floor and the reported power
is the sum of signal power and noise power.
\end{solutionorlines}

\end{questions}

\vspace{1.0em}
\noindent\textbf{Documentation:} \rule[-2pt]{0.6\linewidth}{0.4pt}

\end{document}
